\documentclass[11pt]{article}
\usepackage[utf8]{inputenc}
\usepackage[margin=1in]{geometry}
\usepackage{amsmath}
\usepackage{hyperref}
\usepackage{booktabs}
\usepackage{xcolor}

\title{\textbf{Vehicle Price Prediction: Model Optimization and Framework Analysis}}
\author{LightGBM usage and advantages}
\date{May 2026}

\begin{document}

\maketitle

\section{Executive Summary}
This report outlines the development and optimization of a high-accuracy vehicle price prediction model using \textbf{LightGBM}. Through a combination of logarithmic target transformation, strategic data cleaning, and rigorous hyper-parameter tuning, the model achieved a final \textbf{Mean Absolute Percentage Error (MAPE) of 12.10\%} and a \textbf{Mean Absolute Error (MAE) of \$1,612.18}.

\section{The LightGBM Framework}

\subsection{Definition and Mechanics}
LightGBM (Light Gradient Boosting Machine) is a high-performance, gradient-boosting framework that uses tree-based learning algorithms. Unlike standard decision trees that grow level-by-level, LightGBM grows trees \textbf{leaf-wise}. It chooses the leaf that results in the greatest decrease in loss, allowing for more complex patterns to be captured with fewer total nodes.



\subsection{Suitability for Vehicle Data}
LightGBM is uniquely suited for the "Used Vehicles" dataset for several reasons:
\begin{itemize}
    \item \textbf{Native Categorical Support:} It handles features like \texttt{manufacturer} and \texttt{model} without the need for extensive one-hot encoding, preserving the relationship between categories.
    \item \textbf{Histogram-based Learning:} By binning continuous features like \texttt{odometer} into discrete intervals, it significantly reduces training time and memory usage while maintaining high accuracy.
    \item \textbf{Robustness to Outliers:} Gradient boosting naturally focuses on "hard" samples (residuals), which helps the model learn the nuanced price differences in rare car configurations.
\end{itemize}

\section{Notebook Structure and Workflow}
The project was structured into a logical pipeline designed to move from raw data to a refined valuation engine:
\begin{enumerate}
    \item \textbf{Feature Engineering:} Introduction of \texttt{car\_age} to replace \texttt{year}, providing a more linear representation of value decay.
    \item \textbf{Target Transformation:} Conversion of prices to log-space to optimize for relative (percentage) error rather than absolute dollar error.
    \item \textbf{Data Filtration:} Removal of "noise" rows (scam listings under \$2,000 and luxury outliers over \$85,000) to ensure the model learned authentic market trends.
    \item \textbf{Hyper-parameter Optimization:} Implementing a slow learning rate (0.01) and deep trees (124 leaves) over 10,000 iterations to capture the subtle nuances of the national vehicle market.
\end{enumerate}

\section{Methodology: The Logarithmic Advantage}
Vehicle depreciation is fundamentally a \textbf{multiplicative process} (e.g., a car loses 15\% of its value annually). Predicting raw price often leads to high errors in lower price brackets. By using $y' = \ln(1 + \text{price})$, we:
\begin{itemize}
    \item Symmetrize the distribution of the target variable.
    \item Ensure that a 10\% error on a \$5,000 car is penalized just as heavily as a 10\% error on a \$50,000 car.
\end{itemize}



\section{Results and Performance}
The following table summarizes the evolution of the model performance:

\begin{table}[h]
\centering
\begin{tabular}{@{}lll@{}}
\toprule
\textbf{Model Stage} & \textbf{MAE} & \textbf{MAPE (\%)} \\ \midrule
Baseline (Raw Price) & \$1,757.51 & - \\
Intermediate (Log + Basic Params) & \$1,879.84 & 15.86\% \\
\textbf{Final Optimized Model} & \textbf{\$1,612.18} & \textbf{12.10\%} \\ \bottomrule
\end{tabular}
\caption{Comparison of model performance across optimization stages.}
\end{table}

\section{Conclusion}
The final model demonstrates that predictive accuracy in the used car market depends heavily on data quality and the mathematical alignment of the loss function. By utilizing LightGBM's advanced tree-growth strategies and transforming the target variable, we have developed a model that effectively filters market noise to provide reliable valuations.

\end{document}